% WHAT A RULE BETWEEN AN ALIGNMENT'S ROWS TAKES.
%
% A rule \noalign leaves in an alignment's own list takes whatever it does
% not settle for itself FROM THE BOX THE PREAMBLE IS PACKED AS. For a
% \halign that box is an hbox as wide as the alignment, so a running width
% becomes the alignment's width -- which is how \hline spans a table. For a
% \valign it is a vbox as TALL as the alignment and of no depth, so a
% \vrule with nothing but its width settled comes out
%
%   \rule(10.0+0.0)x0.4      in a \valign 10pt tall
%   \rule(40.0+0.0)x0.4      in \valign to 40pt
%
% This engine settled the width and left the other two running. trip line
% 182 writes \noalign{ \vrule width0pt{ }} in a \valign.
\catcode`\{=1 \catcode`\}=2 \catcode`\#=6 \catcode`\&=4
\tracingonline=1 \showboxdepth=9 \showboxbreadth=9 \hsize=100pt
\font\rip=trip \rip
\message{[A a valign, a rule with every dimension running]}
\setbox0=\hbox{\valign{#\cr\noalign{\vrule}\hbox to 30pt{\vrule height 7pt depth 3pt width 0pt}\cr}}
\showbox0
\message{[B a valign to a height]}
\setbox0=\hbox{\valign to 40pt{#\cr\noalign{\vrule}\hbox to 30pt{\vrule height 7pt depth 3pt width 0pt}\cr}}
\showbox0
\message{[C two columns]}
\setbox0=\hbox{\valign{#\cr\noalign{\vrule}\hbox to 30pt{\vrule height 7pt depth 3pt width 0pt}\cr
  \hbox to 50pt{\vrule height 1pt depth 2pt width 0pt}\cr}}
\showbox0
\message{[done]}
\end
